\section{Motivation and Context}
With more uncrewed aerial vehicles (UAVs) in the air each year, and a growing desire for them to operate in controlled airspaces, there is a concern regarding their ability to maintain separation from hazards in the air. United States and European federal agencies are developing policy that necessitates detect-and-avoid (DAA) systems for UAVs, which parallel the visual and instrument flight rules (VFR \& IFR) that dictate how pilots avoid traffic \cite{easa_easy_access_uas_2024} \cite{rtca_sc228_2025} \cite{serres_wu_2021_nasa_tm}.

While some proposed DAA systems utilise optical and infrared sensors \cite{serres_wu_2021_nasa_tm}, radar sensors may also be suitable, as they can perform sensing over great distances in all weather conditions, due to their relatively low atmospheric attenuation (compared with optical systems) \cite{borden_radar_imaging_airborne_targets_1999}. The inconvenient side effect of their long wavelength property is that target identification by radar systems is less effective than imaging systems that use shorter wavelengths, such as optical trackers. If this challenge is overcome, the implementation of radar as the primary sensing mode or an additional sensing mode will improve the awareness, and thus the safety of UAVs.

In contrast to larger crewed aircraft that can afford great computational resources and large aperture radar, radar systems fit for the size, weight and power (SWaP) conditions of small to medium UAVs may encounter constraints such as low processing power, and a poor signal-to-noise ratio (SNR). Despite the less than ideal conditions for radar on UAVs, they are the fastest growing market for radar systems, projected to grow with a compound annual growth rate (CAGR) of 11.22\% globally through 2030 \cite{mordor_airborne_radars_market_2025}. To make these emerging radar systems more performant, accessible, and safe, it would be valuable to find the most suitable radar signal processing pipeline for UAV DAA systems, revisiting the old problem of radar target tracking in the face of unique hardware conditions.

Software-defined architectures are increasingly preferred over hardware defined radar systems for UAV applications as software modifications can be shipped faster than hardware modifications \cite{mordor_airborne_radars_market_2025}. As a result, graphics processing units (GPUs) are favoured over non-programmable digital signal processors such as application specific integrated circuits (ASICs). GPU-accelerated implementations are also favourable for computationally intensive stages of the radar processing chain, as many core operations, such as fast Fourier transforms (FFTs) and matrix operations, are highly data-parallel and can leverage mature GPU software ecosystems and optimised libraries \cite{zaknoun2024_gpu_radar}. Additionally, the incorporation of AI into radar tracking algorithms has further encouraged these highly parallel, GPU-based pipelines \cite{electronics13214251}. Key open questions remain around balancing the compute saved through GPU parallelism against the cost of data movement and synchronisation, and how unified (shared) CPU--GPU memory on embedded platforms changes this trade-off. 

\todo{Write industry context section heading/intro. Most detail will go in background section so just give high level}

\todo{Current UAV Laws. Australia, USA, Europe, Canada. DAA and how its necessary for UAVs to be integrated into public airspaces.}

\todo{What size vehicles are currently given radar systems. If small UAVs have them, how does the size of the UAV affect the quality/performance of the system.}